\documentclass[11pt,a4paper]{article}
\usepackage{enumitem}

% Remove section numbering
\setcounter{secnumdepth}{0}

% Define document metadata in the preamble
\title{\bfseries Why I use \LaTeX{} and PDF}
\author{Souvik Sarkar \\ \texttt{sounix000@gmail.com}} % Email on a new line as plain text
\date{July 18, 2024} 

\begin{document}

\maketitle % Generates the title, author, email, and date block

It might come across as peculiar that I choose to write my articles in \LaTeX{} and put them out on the web as PDF files, especially when lots of easier markups and browser-friendly alternatives are available. Although I don't owe an explanation to anyone about the choices I make for my personal webpage, I must remind myself of the reasons behind the choice when self-doubts and temptations of easier formats will arise.

\begin{itemize}
\item It looks beautiful and is an archival format. Just for sheer aesthetic pleasure, I will use \LaTeX{} and compile documents to PDF.
\item My writings contain a lot of mathematics, figures/diagrams, code, tables, charts, etc. Nothing beats \LaTeX{} for typesetting and controlling the output in such documents.
\item Mathjax and Katex might be sufficient for mathematical content in HTML, but I just don't like HTML as an output format for such kind of documents. The lack of predictable arrangements of mathematical elements in the output irks me.
\item In my job I write a lot of HTML-friendly content using easier markups that are excellent for code and tables, and occasional illustrations. So for my personal writings I prefer something different.
\item I romanticize academia, especially mathematics. \LaTeX{} gives me a taste of that in my personal writings. If not for anything else, I would continue using \LaTeX{} because of it's association with academia.
\end{itemize}

\end{document}
